% ======================================================================
\chapter{Wireless, cellular, satellite and IoT networks}
\label{ch:wireless}
% ======================================================================

Wireless links break the assumptions behind loss- and delay-based congestion control in several ways at once. Losses come from fading and interference as well as from queue overflow. Capacity on a cellular link can change by an order of magnitude within a second, and the base station's deep per-user buffers turn any overshoot into seconds of delay~\cite{winstein2013sprout}. Wi-Fi aggregation and scheduling make delay noisy. LEO satellite paths see abrupt changes in RTT and capacity, partly caused by frequent satellite handovers~\cite{michel2022starlink}. Constrained IoT devices have little memory or energy for a large model. Learned controllers are attractive here because an agent can learn the link's pattern instead of a fixed rule, and wireless systems often expose \emph{cross-layer} information (signal quality, scheduling, handover notices) that a learned policy can use.

This chapter covers \nChWirelessN{} studies: \nSetCell{} on wireless and cellular links, \nSetSat{} on satellite and space networks, and \nSetIot{} on IoT, sensor and ad hoc networks. Of these, \nChWirelessDesign{} are design studies. The evaluation evidence is weaker than for the Internet. Half of the design studies (\nChWirelessSim) stop at simulation, only \nChWirelessReal{} reach a real network, and \nChWirelessCode{} release code (\cref{sec:wireless-summary}).

\paragraph{Scope note: vehicular ``congestion control''.} In vehicular networking, \emph{decentralized congestion control} (DCC) refers to controlling the load that periodic safety beacons place on the shared radio channel, by adapting each vehicle's beacon rate, transmit power or data rate. Many RL studies address it, and our query retrieves them because they use the phrase ``congestion control''. DCC is a medium-access problem with different objectives (channel busy ratio, awareness range) and no end-to-end flows. We therefore excluded it (EC2, \cref{sec:method}) and include vehicular studies only when they control a transport flow, as in \cref{sec:hsr}.

% ----------------------------------------------------------------------
\section{Cellular and mobile links}
\label{sec:cellular}

\paragraph{Learning on top of existing TCP.} The most influential cellular design is a hybrid. Abbasloo et al.'s DeepCC~\cite{abbasloo2021deepcc} is a DDPG \emph{plug-in} that caps the congestion window of an existing TCP (CUBIC, BBR, Westwood, Illinois) so that the flow meets an application's target delay. It beat both cellular-specific schemes and clean-slate learned controllers, and it shares its plug-in philosophy with Orca, by the same authors (\cref{sec:hybrid}). Later designs tune BBR rather than replace it. BBR-FIT doubles BBR's link utilization on 5G without extra delay by letting an agent set BBR's gain~\cite{xie2022bbr}, and in MIQA a learned bandwidth predictor steers TCP and QUIC controllers on a mmWave testbed~\cite{wu2024miqa}. A digital twin supplies an actor-critic agent with a model of the 5G network for tuning QUIC~\cite{yang2025digital}. Muses~\cite{zhong2022muses} takes a different route to deployability. It imitation-learns an LSTM controller and then distils it into a small decision tree for each environment, which is cheap to run per flow.

\paragraph{Using radio information.} DRL-3R~\cite{chen2021ran} feeds radio-access-network information into the state and redistributes delayed reward to the actions that caused it. Flecto~\cite{serra2025flecto} runs a DQN-family agent inside QUIC on a 5G testbed and reports about 6\% higher throughput than CUBIC with lower RTT (21.8 vs.\ 27.6\,ms). A cross-layer mmWave controller first classifies line-of-sight versus non-line-of-sight conditions~\cite{pan2025intelligent}. At the network side, RLRBM sizes the RAN buffer itself~\cite{ma2022rlrbm}. Radio fronthaul is a special case: Martins et al.\ train TD3 for C-RAN fronthaul congestion and then distil the policy into depth-limited decision trees with sub-microsecond inference~\cite{martins2021policy}, or into closed-form expressions by symbolic regression that match the RL policy's utilization, delay and fairness~\cite{martins2024closed}. Few learned controllers in the corpus are this fast, which is what a hardware datapath requires.

\paragraph{The training pipeline.} Two cellular studies examine the learning pipeline itself. Naqvi et al.'s ablation~\cite{naqvi2022ablation} finds that clipping and scaling of inputs help, that initialization affects results, and that \emph{a higher training reward does not imply better network performance}. Network-segmented model specialization trains separate models for segments of the condition space and raises the QoS of existing DRL controllers~\cite{yao2025unified}. Mobile real-time media is covered in \cref{sec:rtc}. OnRL~\cite{zhang2020onrl}, the largest production evaluation in this setting, belongs there, and so do the Arsenal evaluation platform~\cite{zhang2021arsenal}, the Gymir5G simulation platform~\cite{smirnov2024gymir5g} and 5G WebRTC uplink controllers~\cite{smirnov2024real,liu2025research}.

% ----------------------------------------------------------------------
\section{High-speed rail and vehicular transport}
\label{sec:hsr}

Trains and vehicles combine high mobility with frequent handovers. On high-speed rail, DRL controllers trained on field-measured conditions beat traditional TCP in simulation~\cite{cui2020improving}, a handover-aware design improves goodput~\cite{xu2020handover}, and a space--ground multipath design keeps throughput through frequent handovers~\cite{xu2021deep}. For V2X transport flows, ICCDRL improves throughput, convergence, friendliness and fairness in ns-3~\cite{wang2022intelligent}, and DARA forecasts per-path state with a Transformer and lets a DQN split traffic across vehicular paths~\cite{maglione2026deep}. Apart from DARA, which uses emulation, these studies report simulation results only or do not report their environment.

% ----------------------------------------------------------------------
\section{Wi-Fi, UAV and other wireless links}
\label{sec:wifi}

\paragraph{Loss differentiation.} A classic wireless problem is that a loss-based controller cannot tell a corruption loss from a congestion loss. Several RL designs learn the distinction from features of the loss event and then choose the window response, from an early design that learns wireless versus congestion loss~\cite{rani2019reinforcement} to actor-critic and DRL controllers that report 16--40\% throughput gains~\cite{pradeep2021temporal,midhula2023adaptive}. Others combine DRL with fast recovery for high-loss links~\cite{cao2023drlfcc}, or let Q-learning choose among linear, polynomial and exponential window growth~\cite{molia2024reinforcement}.

\paragraph{Wi-Fi.} Wi-Fi studies are recent and often address fairness. A DQN switches each flow among CUBIC, BBR and PCC to prevent extreme unfairness and starvation~\cite{shrestha2024fairness}. A TD3 agent at the home router throttles bulk torrent traffic to protect interactive flows~\cite{shawish2025wifi}. A DRL regulator tunes BBR's pacing and window gains on a Wi-Fi~7 campus network using eBPF signals, reporting 22\% more throughput than BBRv3 and a Jain index of 0.93~\cite{li2026sac}. Nuwa tunes a receiver-side heuristic with PPO to limit throughput loss when a device switches access points~\cite{gong2022nuwa}. Two designs move the agent off the device. Jazz and TCCO run a Transformer-based DRL agent on an edge server that controls the window of nearby Wi-Fi clients, and both were evaluated on real Wi-Fi testbeds~\cite{wang2025edge,zhang2026transformer}.

\paragraph{UAV and contested links.} Flying ad hoc networks have capacity that swings with distance and attitude. PPO with a continuous action keeps utilization high and p95 RTT low in ns-3~\cite{cai2026deep}, and an A2C controller for UAV-assisted wireless networks improves on a Q-learning baseline in simulation~\cite{feng2023deep}. Under jamming, a hierarchical DRL agent jointly chooses physical-layer anti-jamming actions and the congestion response~\cite{yuan2024drl}. Multi-band links~\cite{meng2024multi} and private wireless networks with in-band telemetry~\cite{bhat2024multi} have also been studied.

% ----------------------------------------------------------------------
\section{Multipath over heterogeneous radios}
\label{sec:wireless-multipath}

Phones and vehicles with Wi-Fi and cellular interfaces are the main use case for multipath transport. SmartCC~\cite{li2019smartcc} applies asynchronous Q-learning with hierarchical tile coding to MPTCP. He et al.'s DeepCC~\cite{he2021deepcc} uses multi-agent DRL with self-attention and converges about 50\% faster than prior DRL designs. RMTC adds pacing and considers TCP-friendliness and head-of-line blocking~\cite{he2023learning}, and a two-stage design targets AR/VR traffic~\cite{xu2021reinforcement}. Leopard~\cite{liu2023leopard} learns offline and then fine-tunes online. In new conditions it converges 20\% faster than MPCC and achieves 13\% more throughput than DRL-CC on a real testbed. Multipath designs for satellite and IoT links are covered below.

% ----------------------------------------------------------------------
\section{Satellite and space networks}
\label{sec:satellite}

Satellite links combine long and changing RTTs, non-congestion loss and, for LEO constellations, frequent handovers. AUTO~\cite{li2021auto} uses multi-objective RL with a preference-adaptation model for satellite--ground links. It beats the state of the art in emulation and on the real Internet, is robust to random loss and rapid changes, and stays fair to other algorithms. LEO-specific designs make the agent aware of handovers. Creo uses DRL while a connection is stable and cross-layer handover notifications around a handover, cutting handover recovery time by 60.5\% in simulation~\cite{yan2025handover}. Other designs target hybrid HEO--LEO constellations~\cite{liu2025intelligent} or integrated space--air--ground networks with heterogeneous multi-agent~\cite{gui2025multi} or hierarchical~\cite{lv2026qos} agents. For non-terrestrial QUIC, RL rate control combined with streaming FEC shortens in-order delivery delay on lossy long-delay links~\cite{yu2023low}. Deep-space and delay-tolerant networks are an extreme case where end-to-end feedback barely exists. Here, RL has been combined with an incomplete-information game to cut bundle aborts in satellite IoT~\cite{wang2019reinforcement}, applied to TCP-like window control with an optimistic actor-critic~\cite{masood2022intelligent}, and used for multipath deep-space video~\cite{ha2023intelligent}.

\paragraph{Simulation versus emulation.} Valentine et al.~\cite{valentine2026emulation} ran the same controllers (CUBIC, BBR, LEO-specific designs and an RL controller) in an OMNeT++/INET simulator and in emulation with real stacks. The simulator misses stack effects such as pacing, SACK/RACK recovery and kernel timing, and these effects change the RL controller's behavior. They conclude that emulation is necessary to validate learned controllers for LEO, which matters because most satellite RL studies in the corpus report simulation results only or do not report their environment (\cref{fig:evaluation}).

\finding{In wireless, satellite and IoT settings, half of the design studies are evaluated in simulation only, yet simulators of these links omit stack effects that change how learned controllers behave. Results from this part of the literature should be treated as provisional until they are reproduced in emulation or on hardware.}

% ----------------------------------------------------------------------
\section{IoT, sensor and ad hoc networks}
\label{sec:iot}

\paragraph{Constrained application protocols.} CoAP, the IoT's REST protocol over UDP, has its own congestion control based on retransmission timeouts. RL agents adapt its parameters~\cite{jayaudhaya2023acoco,donta2023icocoa}, choose how many copies of a message a node sends~\cite{tariq2023ai}, or jointly control CoAP and QUIC with traffic-class-aware rewards that keep critical traffic below 100\,ms~\cite{choi2026energy}. An MQTT broker can also be controlled with DDPG~\cite{janani2026adaptive}. One of the two earliest studies in the corpus, a Q-learning TCP for memory-constrained IoT devices~\cite{li2016learning}, already used function approximation to reduce memory use. Limited resources remain the main concern in this setting. KONTROL offloads the DRL agents of light clients to a central congestion-control server~\cite{wang2025kontrol}, and SymDistill-CC distils PPO into symbolic decision trees for power-IoT terminals~\cite{lu2026wapi}.

\paragraph{Sensor, mesh and ad hoc networks.} In wireless sensor networks the agent often runs at intermediate nodes and controls the forwarding rate or buffer. Examples include an energy-aware Q-learning controller~\cite{trinh2018energy}, a learning-automata controller for Zigbee that reports lower energy and delay~\cite{zhou2023congestion}, cache-aware DRL controllers~\cite{alipio2024cache,mishra2026explainable}, and DLCC for 6LoWPAN, which formulates congestion control as a non-cooperative Markov game with a Lagrangian loss constraint~\cite{hou2023deep}. Hop-by-hop RL flow control appears in lossy and maritime networks~\cite{forero2022congestion,forero2023learning}, and RL-based TCP variants appear in MANETs, mesh and ad hoc networks~\cite{molia2023tcp,mahajan2025reinforcement,li2026lightweight,vazaios2026multi,andradezambrano2025performance}. Delay-tolerant networking motivated one of the first RL congestion controllers in the corpus, Smart-DTN-CC~\cite{silva2019congestion}.

\paragraph{Smart grid and industrial IoT.} Power-grid and industrial IoT communication is a growing niche, with seven studies covering smart meters, power IoT, substations and industrial multipath transport~\cite{andradezambrano2024reinforcement,li2022noninvasive,wang2024research,wang2025transal,lu2026wapi,abedini2026digital,an2026mptcp}. Its distinguishing requirements are reliability and security. One study, for example, couples congestion control with recovery from GOOSE-message flooding attacks~\cite{abedini2026digital}. Two IoT studies evaluate beyond simulation. TransAL-CC evaluates in Pantheon emulation~\cite{wang2025transal}. Kutsevol et al.\ train a \emph{semantic} congestion controller for cyber-physical systems, whose objective is the accuracy of a remote estimator rather than throughput, offline on real traces and validate it on hardware~\cite{kutsevol2025learning}.

% ----------------------------------------------------------------------
\section{Summary}
\label{sec:wireless-summary}

\begin{itemize}
  \item \textbf{Design (\RQ{1}).} Wireless designs act on the window more often than the rate (\nChWirelessActCwnd{} vs.\ \nChWirelessActRate{} design studies). Cross-layer radio state appears almost only in wireless designs, and wireless designs more often place the agent inside the network or at an edge server (\nChWirelessNetwork{} network-side and \nChWirelessCentral{} central agents among \nChWirelessDesign{} design studies). Hybrid designs are less common (\nChWirelessPctHybrid\%) than on the Internet (\nChInternetPctHybrid\%), even though the strongest cellular results come from hybrids.
  \item \textbf{Evidence (\RQ{2}).} \nChWirelessPctSim\% of design studies are evaluated in simulation only and \nChWirelessPctNR\% do not report their environment in the material we could access. Only \nChWirelessReal{} reach a real network and \nChWirelessCode{} release code. Simulators of radio links and satellite stacks miss effects that change learned controllers' behavior~\cite{valentine2026emulation}, so simulation-only results are less dependable here than in the other settings.
  \item \textbf{Deployment (\RQ{3}).} Deployability work focuses on cost. Distillation into trees or closed forms~\cite{zhong2022muses,martins2021policy,martins2024closed,lu2026wapi} and offloading to edge servers~\cite{wang2025edge,zhang2026transformer,wang2025kontrol} are the main techniques. Production evidence is limited to mobile real-time media (OnRL).
\end{itemize}
